\documentclass[11pt]{exam}
\usepackage[utf8]{inputenc}
\usepackage[a4paper,top=2cm,left=2cm,right=2cm,bottom=2cm]{geometry}
\usepackage{import}
\usepackage[T1]{fontenc}
\usepackage[german,ngerman]{babel}
\usepackage[table]{xcolor}
\usepackage{amsthm, esvect}
\usepackage{mathtools}
\RequirePackage{amssymb, amsfonts, amsmath, latexsym, verbatim, xspace}
\usepackage{systeme}
\usepackage{multicol}
\usepackage{multirow}
\usepackage{framed}
\usepackage{array}
\usepackage{ragged2e}
\usepackage{graphicx}
\usepackage{pdflscape}
\usepackage{epstopdf}
\usepackage{booktabs}
\usepackage{pdfpages}
\usepackage{float} % Allows putting an [H] in \begin{figure} to specify the exact location of the figure
\usepackage{wrapfig} % Allows in-line images such as the example fish picture
\usepackage{tikz, graphicx} % tikz
\usepackage[position=top]{subfig}
\usepackage[bookmarks=true]{hyperref}%Links
\usepackage{enumerate}
\usepackage{enumitem}
\usepackage[german,ruled,vlined,linesnumbered,algosection]{algorithm2e}
\usepackage[labelfont=sf]{caption}
\usepackage{lmodern}
\usepackage{lscape}
\usepackage{tabularx}
\usepackage{systeme}
\usepackage{hhline}
\usepackage{subfiles}
\usepackage{stackengine}
\usepackage{setspace}
\usepackage{MnSymbol}
\usepackage{tcolorbox}
\usepackage{bbm}
\usepackage{url}
\usepackage{xparse}
\usepackage{sectsty}
\usepackage{array}
\usepackage{ragged2e}
\usepackage{listings} %Stellt Code-Snippets sauber dar

% definiert die Sprache
\selectlanguage{German}

% add command for different dates in Swiss fashion
\newcommand{\monthwordlong}[1]{\ifcase#1 \or Januar\or Februar\or März\or April\or Mai\or Juni\or Juli\or August\or September\or Oktober\or November\or Dezember\fi}
\newcommand{\monthwordshort}[1]{\ifcase#1\or Jan\or Feb\or Mar\or Apr\or Mai\or Jun\or Jul\or Aug\or Sept\or Okt\or Nov\or Dez\fi}
\newcommand{\datelong}{\number\day.\:\monthwordlong{\the\month}\:\number\year}
\newcommand{\dateshort}{\number\day.\:\monthwordshort{\the\month}\:\number\year}

% define color of links inside a text
\definecolor{linkblue}{RGB}{0, 0, 80}
\definecolor{plotblue}{RGB}{100, 100, 255}
\definecolor{myblue}{RGB}{8,164,214}
\definecolor{myred}{RGB}{143,42,41}
\definecolor{forestgreen}{RGB}{34,139,34}
\definecolor{highdive}{RGB}{10,169,194}
\definecolor{charcoal}{RGB}{74,74,74}
\hypersetup{
	colorlinks=true,
	linkcolor=linkblue,
	filecolor=magenta,
	urlcolor=plotblue,
	citecolor=linkblue,
	pdfpagemode=FullScreen,
}


\lstdefinestyle{mystyle}{ % definiert den style des Code-Snippets
	%backgroundcolor=\color{backcolour},
	%commentstyle=\color{codegreen},
	%keywordstyle=\color{magenta},
	%numberstyle=\tiny\color{codegray},
	%stringstyle=\color{codepurple},
	basicstyle=\ttfamily\footnotesize,
	breakatwhitespace=false,
	breaklines=true,
	captionpos=b,
	keepspaces=true,
	%numbers=left,
	%numbersep=5pt,
	showspaces=false,
	showstringspaces=false,
	showtabs=false,
	tabsize=4
}
\lstset{style=mystyle}
\qformat{\Large\textbf{Aufgabe \thequestion\hfill\thepoints}}
\newlist{partes}{enumerate}{3}
\setlist[partes]{label=(\alph*)}
\newcommand{\parte}{\item}
\newlist{subpartes}{enumerate}{3}
\setlist[subpartes]{label=\roman*)}
\newcommand{\subparte}{\item}
\setlist[enumerate,1]{%(
	leftmargin=*, itemsep=12pt, label={\textbf{\arabic*.)}}}
\newcommand\answerbox{%%
	\fbox{\rule{2in}{0pt}\rule[-0.1ex]{0pt}{4ex}}}
\pointpoints{Punkt}{Punkte}
\hqword{Frage:}
\hpword{Mögliche Punkte:}
\hsword{Erreicht:}
\htword{Total}
\renewcommand*\half{.5}
\renewcommand{\solutiontitle}{\noindent\textbf{Lösung:}\par\noindent}

% Graphed boxes
\newcommand{\karos}[2]{
  \begin{tikzpicture}
    \draw[step=0.4cm,color=cyan,dotted] (0,0) grid (#1 cm ,#2 cm);
  %  \draw((#1)/2 cm,0)--((#1)/2 cm,/#2)/2cm)
  \end{tikzpicture}
}
\usetikzlibrary{arrows, petri, topaths, graphs}
\linespread{1.2} % Line spacing

%=========================================================
\newenvironment{parcols}[1][]{
    \begin{parts}\begin{multicols}{#1}
}{
    \end{multicols}\end{parts}
}
%=========================================================
\newenvironment{itemcols}[1][]{
	\begin{itemize}\begin{multicols}{#1}
		}{
	\end{multicols}\end{itemize}
}

%=========================================================
\newcommand{\antwort}[1]{
	{\setlength{\textwidth}{\textwidth}\fillwithgrid{#1cm}\vspace{5pt}}
}

%=========================================================
% Karriertes Papier für Antwort MIT OVERLAY
\tcbuselibrary{breakable,skins}
\usetikzlibrary{mindmap,trees}
\newlength\tindent
\setlength{\tindent}{\parindent}
\setlength{\parindent}{0pt}
\setlength{\parskip}{0pt}

\newtcolorbox{gridspace}[1]{%
	enhanced,
	breakable,
	blanker,
	boxrule=0pt,
	left=5mm,
	right=5mm,
	top=5.75mm,
	bottom=0mm,
	boxsep=0pt,
	height=#1cm,
	width=16cm,
	before upper={
		\begingroup
		\setlength{\baselineskip}{5.75mm}
		\setlength{\parskip}{0pt}
		\setlength{\lineskip}{0pt}
		% keep every paragraph on the grid rhythm
		%\everypar{\vspace{0pt}\strut}
		% math formulas shifted by one grid row
		\let\oldbracketopen\[
		\renewcommand{\[}{\vspace{-5mm}\oldbracketopen}
		% Align displayed formulas to the grid
		\setlength{\abovedisplayskip}{0pt}
		\setlength{\belowdisplayskip}{0pt}
		\setlength{\abovedisplayshortskip}{0pt}
		\setlength{\belowdisplayshortskip}{0pt}
		% Enumerate follows the same rhythm
		\setlist[enumerate]{leftmargin=6mm,itemsep=0mm,topsep=0pt,parsep=0pt,partopsep=0pt}
		\setlist[itemize]{leftmargin=6mm,itemsep=10mm,topsep=0pt,parsep=0pt,partopsep=0pt}},
	after upper={
		\endgroup},
	underlay={
		\draw[step=5mm,color=gray!55] (interior.south west) grid (interior.north east);}
}

\pagestyle{headandfoot}
\firstpageheader{}{}{\teacher}
\runningheader{}{}{\teacher}
\firstpagefooter{}{}{Seite \thepage\:von \numpages}
\runningfooter{}{}{Seite \thepage\:von \numpages}
%%
% Math Arrows
%%
\newcommand{\same}{\ensuremath{\Longleftrightarrow}}
\renewcommand{\to}{\ensuremath{\longrightarrow}}
\newcommand{\qsame}{\quad\same\quad}
\newcommand{\dsame}{\ensuremath{:\Leftrightarrow}}
\newcommand{\qdsame}{\quad\dsame\quad}
\newcommand{\ldsame}{\ensuremath{:\Longleftrightarrow}}
\newcommand{\samedef}{\ensuremath{\us {\text{def}} \Leftrightarrow}}
\newcommand{\qsamedef}{\quad\samedef\quad}
\newcommand{\lsamedef}{\ensuremath{\us {\text{def}} \Longleftrightarrow}}
\newcommand{\qlsamedef}{\quad\lsamedef\quad}


%%
% Mengendefinitionen
%%
\newcommand{\M}{\ensuremath{\mathcal{M}}}
\newcommand{\N}{\ensuremath{\mathbb{N}}}
\newcommand{\Z}{\ensuremath{\mathbb{Z}}}
\newcommand{\Q}{\ensuremath{\mathbb{Q}}}
\newcommand{\I}{\ensuremath{\mathbb{I}}}
\newcommand{\R}{\ensuremath{\mathbb{R}}}
\newcommand{\C}{\ensuremath{\mathbb{C}}}
\renewcommand{\S}{\ensuremath{\mathbb{S}}}
\newcommand{\Prim}{\ensuremath{\mathbb{P}}}
\newcommand{\0}{\ensuremath{\emptyset}}
\renewcommand{\Re}{\ensuremath{\operatorname{Re}}}
\renewcommand{\Im}{\ensuremath{\operatorname{Im}}}
\newcommand{\strich}{\ensuremath{\big\vert}}

%%
% Mengenoperation
%%
\renewcommand{\nin}{\not\in}


%%
% Logisches Und und Oder
%%
\newcommand{\sand}{\ensuremath{\; \land \;}}
\newcommand{\sor}{\ensuremath{\; \lor \;}}
\renewcommand{\*}{\ensuremath{\cdot}}
\newcommand{\oder}{\vee}
\newcommand{\n}{\neg}
\newcommand{\und}{\wedge}


%%
% Bei Integral, Summe,... die Grenzen über bzw.
% unter das Zeichen
%%
\renewcommand{\d}[1]{\ensuremath{\operatorname{d}\!{#1}}}
\newcommand{\Sum}{\sum\limits}
\newcommand{\Prod}{\prod\limits}
\newcommand{\Min}{\min\limits}
\newcommand{\Max}{\max\limits}
\newcommand{\Lim}{\lim\limits}
\newcommand{\Inf}{\inf\limits}
\newcommand{\Sup}{\sup\limits}
\newcommand{\Liminf}{\liminf\limits}
\newcommand{\Limsup}{\limsup\limits}
%\renewcommand{\Cup}{\bigcup\limits}
%\renewcommand{\Cap}{\bigcap\limits}
\newcommand{\Otimes}{\bigotimes\limits}
\newcommand{\Oplus}{\bigoplus\limits}


%%
% Arcus Cotangens, Logarithmus
%%
\let\oldsin\sin
\renewcommand{\sin}[1]{\oldsin{\left(#1\right)}}
\let\oldcos\cos
\renewcommand{\cos}[1]{\oldcos{\left(#1\right)}}
\let\oldtan\tan
\renewcommand{\tan}[1]{\oldtan{\left(#1\right)}}
\let\oldarcsin\arcsin
\renewcommand{\arcsin}[1]{\oldarcsin{\left(#1\right)}}
\let\oldarccos\arccos
\renewcommand{\arccos}[1]{\oldarccos{\left(#1\right)}}
\let\oldarctan\arctan
\renewcommand{\arctan}[1]{\oldarctan{\left(#1\right)}}
\let\oldln\ln
\renewcommand{\ln}[1]{\oldln{\left(#1\right)}}
\newcommand{\Log}[2]{\ensuremath{\operatorname{log}_{#1}\left(#2\right)}}
\newcommand{\e}{\operatorname{e}}
\renewcommand{\exp}[1]{\ensuremath{\operatorname{\e}^{#1}}}


%%
% Griechische Buchstaben
%%
\renewcommand{\epsilon}{\ensuremath{\varepsilon}}
\renewcommand{\phi}{\varphi}
\renewcommand{\l}{\lambda}
\renewcommand{\t}{\tau}


%%
% Bezeichnungen
%%
\newcommand{\Rgn}{\R_{>0}} % R grösser Null
\newcommand{\Rad}{\operatorname{Rad}}   % Radikal
\newcommand{\kgV}{\operatorname{kgV}}   % Kleinster gemeinsame Vielfache
\newcommand{\ggT}{\operatorname{ggT}}   % Grösster gemeinsame Teiler
\newcommand{\winkel}{\sphericalangle}          % Winkelsymbol
\DeclarePairedDelimiter\abs{\lvert}{\rvert} % besserer Betrag
\makeatletter
\let\oldabs\abs
\def\abs{\@ifstar{\oldabs}{\oldabs*}}
\makeatother


%%
% Vektorgeometrie
%%
\renewcommand{\v}[1]{\vv{#1}}       % besserer Pfeil über Buchstaben
\renewcommand{\vec}[2]{\begingroup\renewcommand{\arraystretch}{1}\left(\begin{array}{c} #1 \\ #2 \end{array}\right)\endgroup}   % two dimensional vector
\renewcommand{\Vec}[3]{\begingroup\renewcommand{\arraystretch}{1}\left(\begin{array}{c} #1 \\ #2 \\ #3 \end{array}\right)\endgroup}      % three dimensional vector


%%
% Text
%%
\newcommand*{\Ueq}[1]{\mathrel{\mathop{=}\limits_{\text{#1}}}}
\newcommand*{\Oeq}[1]{\mathrel{\stackrel{\text{#1}}{=}}}
\newcommand{\utext}[2]{\underbrace{#1}_{\substack{#2}}}
\newcommand{\otext}[2]{$\overbrace{\textrm{#1}}^{\textrm{#2}}$}
\newcommand{\lineortext}[1]{\hspace{0,1cm}\underline{\hspace{0.1cm}\hphantom{#1}\hspace{0.1cm}}\hspace{0,1cm}} % Lückentext erstellen
\newcommand{\gap}[1]{\rule{#1 cm}{1pt}}
%=========================================================
\newcommand{\theme}{Potenzen \& Wurzeln}
\newcommand{\teacher}{Jonas Guler}
%=========================================================
\begin{document}
\fontsize{14}{14}\selectfont
\begin{center}
    \huge\bf\theme
\end{center}

\begin{questions}
%=========================================================
%-------------------TESTFRAGEN
%=========================================================
\question[4]
Vereinfache den folgenden Ausdruck so weit wie möglich.
\begin{parcols}[2]
    \part $6x^8:\left(8x^2\*3y^4\right)$
    \part $\left(\frac{x}{3}\right)^4\*\left(\frac{27}{x^2}\right)^2$
\end{parcols}

%---------------------------------------------------------
\question[4]
Vereinfache den folgenden Ausdruck so weit wie möglich
\begin{parcols}[2]
    \part $\left(6a^9\*4b^4\right):8a^3$
    \part $\left(3ab^2\right)^3\*\frac{\left(a^2b^3\right)^4}{3a^2b^2}$
\end{parcols}

%---------------------------------------------------------
\question[4]
Betrachte den folgenden Lösungsweg. Kennzeichne, wo einen Fehler gemacht wurde und korrigiere die Umformungen ab diesem Punkt.
\begin{align*}
    \left(3ab^2\right)^3\*\frac{\left(a^2b^3\right)^4}{3a^2b^2}&=27a^4b^5\*\frac{a^6b^7}{3a^2b^2}\\
    &= 27a^4b^5\*\frac{1}{3}a^4b^5\\
    &= 9a^8b^{10}
\end{align*}

%---------------------------------------------------------
\question[3]
Betrachte den folgenden Lösungsweg. Kennzeichne, wo einen Fehler gemacht wurde und korrigiere die Umformung ab diesem Punkt.
\begin{align*}
    \left(\frac{a}{3}\right)^4\*\left(\frac{27b}{a^2}\right)^2 &= \frac{a^4}{3}\*\frac{\left(27b\right)^2}{a^2}\\
    &= \frac{a^4}{3}\*\frac{27b^2}{a^2} \\
    &= 9a^2b^2
\end{align*}






%======================= WURZELN =========================
\question[4]
Schreibe die folgenden Terme ohne Wurzelzeichen.
\begin{parcols}[2]
    \part $\sqrt{(x+1)^2-4x}$ %Alegbra 910 S132 A3e
    \part $\sqrt{\frac{25a^2b^8}{16}}$ %Algebra 910 S132 A3h
\end{parcols}

%---------------------------------------------------------
\question[4]
Schreibe die folgenden Terme ohne Wurzelzeichen.
\begin{parcols}[2]
    \part $\sqrt{(a-2)^2+8a}$
    \part $\sqrt{\frac{49(2x+y)^2}{2401x^2}}$
\end{parcols}



%---------------------------------------------------------
\question[4]
Betrachte den folgenden Lösungsweg. Kennzeichne, wo einen Fehler gemacht wurde und korrigiere die Umformungen ab diesem Punkt.
\begin{normalsize}
    \begin{align*}
        \left(\sqrt{r}-\sqrt{s}\right)^2-\left(\sqrt{r}+\sqrt{s}\right)^2 &= \left(\left(\sqrt{r}\right)^2-2\*\sqrt{r}\sqrt{s}+\left(\sqrt{s}\right)^2\right)-\left(\left(\sqrt{r}\right)^2+2\*\sqrt{r}\sqrt{s}+\left(\sqrt{s}\right)^2\right)\\
        &= \left(r-2\*\sqrt{r}\sqrt{s}+s\right)-\left(r+2\*\sqrt{r}\sqrt{s}+s\right)\\
        &= r-2\*\sqrt{r+s}+s-r-2\*\sqrt{r+s}-s\\
        &=-4\*\sqrt{r+s}
    \end{align*}
\end{normalsize}

%---------------------------------------------------------
\question[3]
Betrachte den folgenden Lösungsweg. Kennzeichne, wo einen Fehler gemacht wurde und korrigiere die Umformungen ab diesem Punkt.
\begin{normalsize}
    \begin{align*}
        \left(\sqrt{x}-5\sqrt{y}\right)\left(\sqrt{x}+5\sqrt{y}\right)-\left(\sqrt{x}-5\sqrt{y}\right)^2 &= \left(\sqrt{x}\right)^2-\left(5\sqrt{y}\right)^2-\left(\left(\sqrt{x}\right)^2-2\*5\*\sqrt{x}\*\sqrt{y}+\left(5\sqrt{y}\right)^2\right) \\
        &= \left(x-5y\right)-\left(x-10\sqrt{x}\sqrt{y}+5y\right) \\
        &= x-5y-x+10\sqrt{x+y}-5y \\
        &= 10\sqrt{x+y}-10y
    \end{align*}
\end{normalsize}


%---------------------------------------------------------
\question[2]
Entscheide in der Tabelle unten, ob die Aussagen wahr oder falsch sind. Die Antworten müssen nicht begründet werden.\newline
\begin{small}
    Bewertung: Für $0-2$ korrekte Antworten gibt es keine Punkte; $3$ korrekte Antworten einen Punkt und falls alle Antworten korrekt sind, gibt es $2$ Punkte.
\end{small}
\begingroup
\renewcommand{\arraystretch}{1.5}
\begin{table}[h]
    \centering
    \begin{tabular}{|p{12cm}|c|c|}\hline
         & wahr & falsch \\\hline
        $\sqrt{a}$ hat zwei Lösungen, wobei eine positiv und eine negativ ist. & & \\\hline
        $\displaystyle\sqrt{9x}:\sqrt{x}+\left(\sqrt{x-2}\right)^2=x+1$ & & \\\hline %Algebra 910 S132 A1c
        In der Gleichung $2^\triangle=8^{3}$ muss das Symbol $\triangle$ mit der Zahl $12$ ersetzt werden, damit die Gleichheit stimmt. & & \\\hline
        In der Gleichung $\square^{12}=16^3$ muss das Symbol $\square$ mit der Zahl $4$ ersetzt werden, damit die Gleichheit stimmt.  & & \\\hline
        $\displaystyle\left(\sqrt{x-2}\right)^2=x+1-\sqrt{9x}:\sqrt{x}$ & & \\\hline %Algebra 910 S132 A1c
        In der Gleichung $3^\triangle=81^{3}$ muss das Symbol $\triangle$ mit der Zahl $12$ ersetzt werden, damit die Gleichheit stimmt. & & \\\hline
        In der Gleichung $\square^{8}=16^4$ muss das Symbol $\square$ mit der Zahl $4$ ersetzt werden, damit die Gleichheit stimmt.  & & \\\hline
    \end{tabular}
\end{table}
\endgroup

\end{questions}
\end{document}